\section{Reverse a Linked List (Iterative)}
\label{sec:ll-reverse-iterative}

\problemstatement{Reverse a singly linked list in place and return the new
head, using $O(1)$ extra space.}

\begin{patternbox}
Reversing a singly linked list is the archetypal \textbf{pointer-rewiring}
loop: walk the list carrying three pointers -- \texttt{prev},
\texttt{curr}, \texttt{next} -- and at each node flip \texttt{curr->next} to
point backward at \texttt{prev}.
\end{patternbox}

\subsection*{Flip each link, remember the rest}

The danger in rewiring is losing the remainder of the list. So before
flipping \texttt{curr->next} to \texttt{prev}, save \texttt{curr->next} in a
temporary \texttt{next}. Then set \texttt{curr->next = prev}, advance
\texttt{prev} to \texttt{curr}, and \texttt{curr} to the saved
\texttt{next}. When \texttt{curr} becomes null, \texttt{prev} is the new
head (the old tail).

\begin{ideabox}[Save Next Before You Overwrite It]
Each step overwrites \texttt{curr->next}, which is the only link to the rest
of the list -- so it must be stashed first. This save--flip--advance trio is
the canonical in-place reversal and recurs inside reverse-in-groups and
palindrome checks.
\end{ideabox}

\subsection*{Algorithm}

\begin{enumerate}
  \item \texttt{prev = null}, \texttt{curr = head}.
  \item While \texttt{curr}: save \texttt{next = curr->next}; set
        \texttt{curr->next = prev}; advance \texttt{prev = curr},
        \texttt{curr = next}.
  \item Return \texttt{prev}.
\end{enumerate}

\subsection*{Dry run}

\dryrun{List $1\to2\to3$.}

\begin{itemize}
  \item prev=null,curr=1: flip $1\to$null; prev=1,curr=2.
  \item flip $2\to1$; prev=2,curr=3. flip $3\to2$; prev=3,curr=null.
  \item New head $3\to2\to1$.
\end{itemize}

\subsection*{C++ implementation}

\begin{lstlisting}
#include <bits/stdc++.h>
using namespace std;

struct ListNode {
    int val; ListNode* next;
    ListNode(int v) : val(v), next(nullptr) {}
};

ListNode* build(const vector<int>& a) {
    ListNode dummy(0), *t = &dummy;
    for (int x : a) { t->next = new ListNode(x); t = t->next; }
    return dummy.next;
}

ListNode* reverse(ListNode* head) {
    ListNode* prev = nullptr;
    ListNode* curr = head;
    while (curr) {
        ListNode* next = curr->next;              // save the rest
        curr->next = prev;                        // flip
        prev = curr;                              // advance
        curr = next;
    }
    return prev;                                  // new head
}

void print(ListNode* h) { for (; h; h = h->next) cout << h->val << " "; cout << "\n"; }

int main() {
    ListNode* head = build({1, 2, 3, 4});
    head = reverse(head);
    print(head);                                  // 4 3 2 1
    return 0;
}
\end{lstlisting}

\begin{complexitybox}
\textbf{Time Complexity.} $O(n)$. \textbf{Space Complexity.} $O(1)$.
\end{complexitybox}

\begin{takeawaybox}
Iterative reversal is the save--flip--advance loop with three pointers,
$O(1)$ space. It is the most important pointer-manipulation pattern in the
chapter -- committed to muscle memory, it makes reverse-in-$k$-groups,
palindrome, and add-two-numbers straightforward.
\end{takeawaybox}
